\subsection*{Grandezas e unidades de medida}
    \subsubsection*{Distância e comprimento}
    Distância e comprimento medem o quão longe dois pontos estão entre si, e no caso de arestas essa medida é
    denotada como comprimento. A unidade de medida base utilizada para comprimento é o metro (símbolo ${m}$), além do
    metro existem seus múltiplos que são igualmente utilizados dependendo da distância/comprimento aferido.
    \begin{table}[H]
        \resizebox{\columnwidth}{!}{%
            \begin{tabular}{l|l|l}
                \rowcolor{gray!25}
                \textbf{Nome} & \textbf{Sigla} & \textbf{Equivalência} \\ \hline
                picometro     & $pm$           & $10^{-12}m$           \\ \rowcolor{gray!10}
                nanometro     & $nm$           & $10^{-9}m$            \\
                micrometro    & $\mu m$        & $10^{-6}m$            \\ \rowcolor{gray!10}
                milimetro     & $mm$           & $10^{-3}m$            \\
                centímetro    & $cm$           & $10^{-2}m$            \\ \rowcolor{gray!10}
                decímetro     & $dm$           & $10^{-1}m$            \\
                metro         & $m$            & $1m$                  \\ \rowcolor{gray!10}
                decâmetro     & $dam$          & $10m$                 \\
                hêctometro    & $hm$           & $100m$                \\ \rowcolor{gray!10}
                quilometro    & $km$           & $1000m$
            \end{tabular}%
        }
        %\caption{Unidades medida de comprimento e distância}
    \end{table}
    Outras medidas de comprimento são: \\
    \textbf{Perímetro:} comprimento do contorno de uma fígura geométrica. \\
    \textbf{Raio:} distância entre o centro de uma circunferência até seu contorno ou superfície.\\
    \textbf{Diâmetro:} comprimento de reta que passe pelo centro da circunferência e cujo seus pontos de
    início e fim estejam sobre a circunferência\\
    \textbf{Circunferência:} perímetro de uma circunferência.\\

    \begin{figure}[H]
        \centering
        \includegraphics[width=0.7\columnwidth]{assets/image_placeholder.png}
        \caption{Raio, diâmetro e circunferência de um círculo}
    \end{figure}

    \subsubsection*{Área}
    Área é a medida que expressa a quantidade de espaço bidimenssional ocupado por uma fígura geométrica,
    área de superfície é o equivalente da área para uma superfície ou face de um objeto tridimenssional. A unidade base
    para área é o metro quadrado (símbolo ${m^2}$). Seus múltiplos também acompanham o termo 'quadrado',
    e a equivalência é a mesma do metro, porém elevada ao quadrado. Ex.: $1m = 10^2cm$ e $1m^2 = (10^2cm)^2 = 10^4cm^2$,
    lembrar essa regra pode facilitar na hora da conversão de área.

    \begin{table}[H]
        \resizebox{\columnwidth}{!}{%
            \begin{tabular}{l|l|l}
                \rowcolor{gray!25}
                \textbf{Nome}       & \textbf{Sigla} & \textbf{Equivalência}       \\ \hline
                picometro quadrado  & $pm^2$         & $10^{-24}m^2$               \\ \rowcolor{gray!10}
                nanometro quadrado  & $nm^2$         & $10^{-18}m^2$               \\
                micrometro quadrado & $\mu m$        & $10^{-12}m^2$               \\ \rowcolor{gray!10}
                milimetro quadrado  & $mm^2$         & $10^{-6}m^2$                \\
                centímetro quadrado & $cm^2$         & $10^{-4}m^2$                \\ \rowcolor{gray!10}
                decímetro quadrado  & $dm^2$         & $10^{-2}m^2$                \\
                metro quadrado      & $m^2$          & $1m^2$                      \\ \rowcolor{gray!10}
                decâmetro quadrado  & $dam^2$        &                             \\ \rowcolor{gray!10}
                are                 & $are$          & \multirow{-2}{*}{$100m^2$}  \\
                hêctometro quadrado & $hm^2$         &                             \\
                hectare             & $ha$           & \multirow{-2}{*}{$10^4m^2$} \\ \rowcolor{gray!10}
                quilometro quadrado & $km^2$         & $10^6m^2$
            \end{tabular}%
        }
        %\caption{Unidades medida de área}
    \end{table}

    \subsubsection*{Volume}
    Assim como a área expressa a quantidade de espaço bidimenssional ocupado por uma fígura geométrica, o volume expressa
    a quantidade de espaço tridimenssional ocupado por um objeto. No caso do volume existe duas unidades de medida base,
    o metro cúbico e o litro, ambos são usados regularmente, na qual o litro é a unidade mais usada para volumes
    pequenos, enquanto o metro cúbico é usado para expressar volumes maiores.

    \begin{table}[H]
        \resizebox{\columnwidth}{!}{%
            \begin{tabular}{l|l|l}
                \rowcolor{gray!25}
                \textbf{Nome}     & \textbf{Sigla} & \textbf{Equivalência}          \\ \hline
                milimetro cúbico  & $mm^3$         & $10^{-9}m^3$                   \\ \rowcolor{gray!10}
                centímetro cúbico & $cm^3$         &                                \\ \rowcolor{gray!10}
                mililitro         & $ml$           & \multirow{-2}{*}{$10^{-6}m^3$}  \\
                decímetro cúbico  & $dm^3$         &                                \\ %\hline
                litro             & $l$            & \multirow{-2}{*}{$10^{-3}m^3$} \\ \rowcolor{gray!10}
                metro cúbico      & $m^3$          & $1m^2$                         \\ %\hline                      
            \end{tabular}%
        }
        %\caption{Unidades medida de volume}
    \end{table}

    \subsubsection*{Ângulo}

    Ângulo é uma medida de inclinação entre duas retas ou dois planos, desde que eles não sejam paralelos entre si, no 
    contexto da geometria os ângulos estão presentes em todos os vértices das fíguras geométricas. 
    As unidades de medida mais comuns para ângulo são graus (símbolo $^{\circ}$), radianos (símbolo $rad$) e gradianos (símbolo $gon$).
    Quanto no escopo das funções trigonométricas a unidade de medida mais utilizada é o radiano, por outro lado, 
    em aplicações de engenharia e na geometria a unidade mais usada é o grau. Considerando um grau $\alpha$ em graus, 
    sua conversão para radianos se dá por $rad(\alpha) = \alpha\cdot\pi/180$ 

    % teste matemática em duas colunas
    % \color{black}
    % \begin{align}
    %     f(x) &= \frac{sen(x)}{cos(x)} \\[10pt]
    %     f(x) &= tan(x)
    % \end{align}
    % \color{mediumgray}

    %\pagebreak
    % \newpage
    % \section*{Ângulos}
    % \lipsum[1]

    % \subsection*{Simetrias de figuras planas}
    % \lipsum[1-2]    